Nous avons appris au chapitre \ref{ch:proc-pascal} à déclarer et à appeler des procédures sans paramètres
nous permettant de répéter plusieurs fois la même séquence de commmandes du
\robot\ à des endroits différents dans le programme. Au chapitre \ref{ch:appel} nous avons
utilisé des commandes du \robot\ qui sont des appels de procédures avec paramètres,
l'usage d'un ou plusieurs paramètres per\-met\-tant à la procédure de faire des séquences
de commandes similaires mais différentes selon la valeur des paramètres. Nous allons maintenant
voir comment déclarer des procédures avec paramètres en \pascal. Commençons
par un exemple simple~:

\bigskip

 {\raggedright 
 \vbox to 0pt{\vspace*{\baselineskip}\vspace*{3cm}
 \includevspacePostScript[-\baselineskip]{CarrePointille}{0.95}{0.95}}%
 \input CarrePointille.tex
 }
La procédure {\idt TraitPointille} permet de tracer un trait en pointillé cons\-ti\-tué par une succession de
{\idt N} pixels noirs séparés par $(${\idt N}$ - 1)$ pixels blancs, le nom\-bre {\idt N} de points étant
passé en paramètre de cette procédure.  La déclaration de la procédure {\idt TraitPointille} indique
que le paramètre {\idt N} est un entier ({\idt integer} en anglais).
La procédure {\idt Carre} permet de dessiner un carré
en pointillé, le nombre {\idt P} de points par côté étant passé en paramètre. L'appel {\idt  Carre$($30$)$}
indique qu'il faut exécuter la procédure {\idt Carre} avec un paramètre {\idt P} égal à 30. De ce fait la
procédure {\idt TraitPointille} est appelée quatre fois avec un paramètre {\idt N} égal à {\idt P} soit
30. On dit que {\idt P} est le {\em paramètre formel}\index{Paramètre formel}  (dont on ne connaît pas encore la valeur
au moment de la déclaration de la procédure et que l'on désigne par un identificateur) et que 
30 est le {\em paramètre effectif}\index{Paramètre effectif} (c'est-à-dire
la valeur du paramètre formel pour cet appel de la procédure).

De manière générale, les déclarations et appels de procédures avec pa\-ra\-mè\-tres ont  la forme
suivante~:

\index{Paramètre formel}\index{Déclaration de procédure}\index{Identificateur}
\begin{center}
\begin{tabular}{@{}|p{12.5cm}|@{}} \hline
{\em Déclaration de procédure avec paramètres~:}\\ \hline \hline
\begin{minipage}[t]{12.5cm}
{\medskip
\mbox{\blc\ \ \ \mcl procedure\blc\hspace{0.3em}\blanc\hspace{0.3em}\idt NomDeLaProcedure%
\ $($ParamFormel1\footnote{\ Un paramètre formel est un identificateur.}\ $:$\ integer$;$}\par 
\mbox{\blc\ \ \ \ \ \ \ \ \ \ \ \ \ \ \ \ \ \ \ \idt \ ParamFormel2\ $:$\ integer\ \dots$);$}\par 
\mbox{\blc\ \ \ \mcl begin}\par 
\mbox{\blc\ \ \ \ \small\rm \ldots\ Instructions du corps de la procédure%
\footnote{\ Les expressions entières figurant dans ces instructions peuvent utiliser les paramètres formels
{\idt ParamFormel1}, \ldots\ de la procédure.}
 (séparées par des `$;$') \ldots}\par 
\mbox{\blc\ \ \ \mcl end\idt$;$}\par 
\medskip}
\end{minipage}\\ \hline 
\end{tabular}
\end{center}

\index{Paramètre effectif}\index{Appel de procédure}
\begin{center}
\begin{tabular}{@{}|p{12.5cm}|@{}} \hline
{\em Appel de procédure avec paramètres~:}\\ \hline \hline
\begin{minipage}[t]{12.5cm}
{\medskip
\mbox{\blc\ \ \ \idt NomDeLaProcedure\ $($ParamEffectif1%
\footnote{\ Un paramètre effectif est une expression entière.}$,$\ ParamEffectif2\ \dots$%
\footnote{\ Les paramètres effectifs sont donnés dans l'ordre où apparaissent les pa\-ra\-mè\-tres
formels correspondants.});$}\par 
\medskip}
\end{minipage}\\ \hline 
\end{tabular}
\end{center}

On remarquera que les paramètres formels figurant entre parenthèses dans la déclaration de
procédure sont séparés par des points-virgules tandis que les paramètres effectifs figurant
entre parenthèses dans l'appel de procédure sont séparés par des virgules. 

Une procédure avec 
paramètres formels peut être appelée plusieurs fois dans un programme avec des 
pa\-ra\-mè\-tres effectifs
qui peuvent être différents à chaque fois.

Les paramètres formels d'une procédure ne sont utilisables que dans le
corps de la procédure. Par conséquent des procédures différentes peuvent avoir des
paramètres formels portant le même nom bien qu'il s'agisse de paramètres différents. 

Ces différents 
points sont illustrés dans l'exemple ci-dessous.

\begin{exemple}[Spirale segmentée]{\rm
~Le programme ci-dessous dessine une spirale constituée de paires de segments de droite
de longueurs successives 1, 2, 3, \ldots, 19~:

 {\raggedright 
 \vbox to 0pt{\vspace*{\baselineskip}\vspace*{4cm}
 \includevspacePostScript[-\baselineskip]{SpiraleSegmentee}{0.95}{0.95}}
 \input SpiraleSegmentee_2.tex
 }
\noindent $\Box$ }\end{exemple}


\noindent \rule{137mm}{0.1mm} %---------------------------------------------
\begin{exercices}{Écrire des programmes \pascal\ comportant des procédures avec
paramètres pour réaliser les dessins suivants~:

\tiret la croix à trois branches transversales, l'hélice coloriée, le mi\-cro-\-or\-di\-na\-teur,
le logo en couleur, la chenille (où le nombre
d'anneaux du corps de la chenille est passé en paramètre de la procédure de dessin), la pirogue 
(où le nombre de rameurs est un paramètre) représentés à la page \pageref{fig:crx3b}$\:$;

\tiret les carrés inscrits les uns dans les autres
ou la version peinte de ce dessin représenté à la figure \ref{fig:carinspt} (page \pageref{fig:carinspt})$\:$;

\tiret l'arabesque inspirée par la coupole du mausolée du Sultan
         Qayt-\-Bay re\-pré\-sen\-tée à la figure \ref{fig:arab2} (page \pageref{fig:arab2})$\:$;

\tiret la péniche de la figure \ref{fig:transports} (page \pageref{fig:transports}) où le nombre
         de hublots est donné en paramètre$\:$;

\tiret le pavage de Golomb de la figure \ref{fig:golomb} (page \pageref{fig:golomb}).

\noindent $\Box$}\end{exercices}
\twocolumn
\begin{figure}[h]
%\hspace{-3.363cm} 
\includevspacePostScript[4.9cm]{Arabesque2}{0.95}{0.95}
 \caption{Arabesque}\label{fig:arab2}
\end{figure}
\vfill
\begin{figure}[h]
%\hspace{-3.363cm} 
\includevspacePostScript[0.0cm]{HeliceEnCouleur}{0.95}{0.95}
 \caption{Hélice coloriée}
 %\label{fig:helice}
\end{figure}
\vfill
\begin{figure}[h]
%\hspace{-3.363cm} 
\includevspacePostScript[0.0cm]{MacintoshSE}{0.95}{0.95}
 \caption{Ancien micro-ordinateur}
 %\label{fig:MacSE}
\end{figure}
\mbox{}
\newpage
\begin{figure}[!h]
%\hspace{-3.363cm} 
\includevspacePostScript[5.3cm]{LogoEnCouleur}{0.95}{0.95}%
\caption{Logo en couleur}%
%\label{fig:Logoc}
\end{figure}
\begin{figure}[!h]
%\hspace{-3.363cm} 
\includevspacePostScript[1.0cm]{Chenille}{0.95}{0.95}%
\caption{Chenille}%
%\label{fig:chenille}
\end{figure}
\begin{figure}[!h]
%\hspace{-3.363cm} 
\includevspacePostScript[1.0cm]{Pirogue_1}{0.95}{0.95}%
\caption{Pirogue}%
%\label{fig:pirog1}
\end{figure}
\vfil
\begin{figure}[!h]
\includevspacePostScript[0.0cm]{Croix3branches}{0.95}{0.95}%
\caption{Croix à trois branches}%
\label{fig:crx3b}%
\end{figure}
\onecolumn
\begin{figure}[h]
\includevspacePostScript[5.7cm]{CarresInscritsPeints}{0.95}{0.95}
 \caption{Motif de carrés inscrits de Horemis}\label{fig:carinspt}
\end{figure}
\begin{corrige}{
Nous donnons un programme, les autres se trouvant sur disquette~:
\input HeliceEnCouleur.tex
\noindent $\Box$}\end{corrige}

\vfill

%%%%%%%%%%%%%%%%%%%%%%%%%%%%%%%%%%%%%%%%%%%%%%%%%%%%%%%%%%%%%%%%%%%%%%%%%%%
\noindent \framebox[\textwidth]{
\begin{minipage}{\myparboxwidth}
\setlength{\parindent}{2em}
\small \sf
\vspace{\abovedisplayskip}
\section*{Addenda}

On distingue en \pascal, les définitions des identificateurs de leurs utilisations. \index{Identificateur}
Les   {\em occurrences de définition} figurent dans une déclaration de
constante, de va\-ria\-ble ou de procédure comme nom de procédure ou de paramètre formel.
Les  {\em occurrences d'utilisation} de ces identificateurs sont les noms de constante, de variable ou de
paramètre formel utilisés dans une expression entière ou les noms de procédure utilisés dans des appels.

 Il se peut qu'un même identificateur ait plusieurs occurrences de définition dans un programme. De ce fait il se
 pose le problème d'établir la correspondance entre les occurrences d'utilisation et de définition d'un même
 identificateur. Par 
exemple l'identificateur peut être utilisé pour nommer une constante du programme et des paramètres
formels dans des procédures. En \pascal, ces procédures peuvent être imbriquées les unes dans les
 autres, comme dans l'exemple ci-après~:
 
\smallskip

\parindent \parindentprog

\mbox{\blc  \idt  \mcl program\idt \label{prog:Np}\glossary{\idt Np}\ Np$;$}\par 
\vspace{\neglskip}\mbox{\blc \blk\idt \mcl uses\idt \ \ldots\ $;$}\par 
\vspace{\neglskip}\mbox{\blc  \blk\idt  \mcl const\idt \ N\ $=$\ 17$;$}\par 

\vspace{\neglskip}\mbox{\blc  \blk\idt  \mcl procedure\idt \ P1\ $($N\ $:$\ integer$)$$;$}\par 

\vspace{\neglskip}\mbox{\blc  \blk\blk\idt  \mcl procedure\idt \ P2\ $($N\ $:$\ integer$)$$;$}\par 
\vspace{\neglskip}\mbox{\blc  \blk\blk\idt  \mcl begin\idt }\par 
\vspace{\neglskip}\vspace{\neglskip}\mbox{\blc  \blk\blk\blk\idt  \cmt  \ \ \ \idt \ \ \ldots}\par 
\vspace{\neglskip}\vspace{\neglskip}\mbox{\blc  \blk\blk\blk\idt  \cmt  \{1\}\idt \ \ \ldots\ N\ \ldots}\par 
\vspace{\neglskip}\vspace{\neglskip}\mbox{\blc  \blk\blk\blk\idt  \cmt  \ \ \ \idt \ \ \ldots}\par 
\vspace{\neglskip}\mbox{\blc  \blk\blk\idt  \mcl end\idt $;$}\par 

\vspace{\neglskip}\mbox{\blc  \blk\idt  \mcl begin\idt }\par 
\vspace{\neglskip}\vspace{\neglskip}\mbox{\blc  \blk\blk\idt  \cmt  \ \ \ \idt \ \ \ldots}\par 
\vspace{\neglskip}\mbox{\blc  \blk\blk\idt  \cmt  \{2\}\idt \ \ P2$($N$)$$;$}\par 
\vspace{\neglskip}\vspace{\neglskip}\mbox{\blc  \blk\blk\idt  \cmt  \ \ \ \idt \ \ \ldots}\par 
\vspace{\neglskip}\mbox{\blc  \blk\idt  \mcl end\idt $;$}\par 

\vspace{\neglskip}\mbox{\blc  \idt  \mcl begin\idt }\par 
\vspace{\neglskip}\vspace{\neglskip}\mbox{\blc  \blk\idt  \cmt  \ \ \ \idt \ \ \ldots}\par 
\vspace{\neglskip}\mbox{\blc  \blk\idt  \cmt  \{3\}\idt \ \ P1$($N$)$$;$}\par 
\vspace{\neglskip}\vspace{\neglskip}\mbox{\blc  \blk\idt  \cmt  \ \ \ \idt \ \ \ldots}\par 
\vspace{\neglskip}\mbox{\blc  \idt  \mcl end\idt $.$}\par 

\vspace{\abovedisplayskip}
\end{minipage}
}
%%%%%%%%%%%%%%%%%%%%%%%%%%%%%%%%%%%%%%%%%%%%%%%%%%%%%%%%%%%%%%%%%%%%%%%%%%%% %%%%%%%%%%%%%%%%%%%%%%%%%%%%%%%%%%%%%%%%%%%%%%%%%%%%%%%%%%%%%%%%%%%%%%%%%%%
%%%%%%%%%%%%%%%%%%%%%%%%%%%%%%%%%%%%%%%%%%%%%%%%%%%%%%%%%%%%%%%%%%%%%%%%%%%
\noindent \framebox[\textwidth]{
\begin{minipage}{\myparboxwidth}
\setlength{\parindent}{2em}
\small \sf
\vspace{\abovedisplayskip}
 
 On dit que le programme a une {\em structure de blocs}. \index{Structure de blocs}Ces blocs sont formés par les procédures imbriquées 
 les unes dans les autres, le programme étant  considéré comme le bloc le plus externe. 
 
 La {\em règle de 
 visibilité}\index{Visibilité} des identificateurs en \pascal\ dit que l'occurrence d'utilisation d'un identificateur
 se réfère à l'occurrence de définition du bloc le plus interne où figure cette occurrence d'utilisation et
 qui contient une occurrence de définition de l'identificateur. 

Dans l'exemple ci-dessus,
\begin{list}{$\scriptscriptstyle\bullet$}{
}
 \item l'occurrence d'utilisation de {\idt N} à la ligne 1 se réfère au paramètre formel de la procédure {\idt P2}$\:$;
\item l'occurrence d'utilisation de {\idt N} à la ligne 2 se réfère au paramètre formel de la procédure {\idt P1}$\:$;
 \item l'occurrence d'utilisation de {\idt N} à la ligne 3 se réfère à la constante 17.
\end{list}

La structure de bloc de \pascal\ est héritée d'{\sc algol 60} \cite{algol60}. Ce langage fut créé par un groupe
d'experts européens et américains  réunis dans le groupe de travail WG2.1 de l'IFIP (International Federation of 
Information Processing) qui regroupe toutes les associations nationales d'informaticiens. C'est l'ancêtre d'une famille
de langages parmi lesquels on compte \pascal, {\sc Ada} et {\sc Modula}.

\vspace{\abovedisplayskip}
\end{minipage}
}
%%%%%%%%%%%%%%%%%%%%%%%%%%%%%%%%%%%%%%%%%%%%%%%%%%%%%%%%%%%%%%%%%%%%%%%%%%%%